\documentclass{article}
\usepackage[utf8]{inputenc}
\usepackage{hyperref}
\usepackage{booktabs}
\usepackage{geometry}
\geometry{a4paper, margin=1in}

\title{Hybrid Knowledge Graph Embeddings for Ranking Security-Event Technique Links: A Small PyKEEN Feasibility Study}
\author{Knowledge Graphs Project Report}
\date{\today}

\begin{document}

\maketitle

\section{Research Question}
Can a compact hybrid knowledge graph, consisting of observed security-event facts and rule-derived technique hypotheses, support link prediction of missing \texttt{process -> suggests\_technique -> ATT\&CK-like technique} triples?

\section{Data and Method}
\subsection{Data Statement}
We constructed a deterministic synthetic enterprise-security knowledge graph with 1,232 labelled triples, 245 entities and 9 relations. The observed-style component contains 944 triples representing pseudonymous users, workstations, servers, processes, file kinds and asset roles. A transparent rule layer derives 288 additional facts. For four synthetic process families, the rule maps a \texttt{has\_family} observation to an ATT\&CK-inspired \texttt{suggests\_technique} hypothesis and stores a provenance relation. The graph labels are semantically grounded in MITRE ATT\&CK technique identifiers, but the events and labels are generated locally rather than collected from real incidents.

\subsection{Method}
We modelled a synthetic enterprise-security knowledge graph as labelled triples. The graph contains users, workstations, servers, processes, files, process families, security roles and ATT\&CK-like techniques. Observed facts include \texttt{authenticates\_to}, \texttt{runs}, \texttt{executed\_by}, \texttt{has\_family}, \texttt{accesses}, \texttt{connects\_to}, and \texttt{has\_role}. A deterministic Datalog-style rule derives \texttt{suggests\_technique} triples from process-family evidence and adds a provenance fact. We randomly split triples 80/10/10 into training, validation and testing sets using seed 42. We trained TransE and ComplEx with a 32-dimensional embedding, Adam optimisation (learning rate 0.01) and 80 epochs. We assessed filtered link prediction using MRR, Hits@1, Hits@3, Hits@10 and mean rank.

\section{Results}
\subsection{Model Comparison}
The following table shows the validation set performance for TransE and ComplEx trained under identical conditions (80 epochs, dimension 32).

\begin{table}[h]
\centering
\begin{tabular}{lrrrrr}
\toprule
\textbf{Model} & \textbf{Filtered MRR} & \textbf{Hits@1} & \textbf{Hits@3} & \textbf{Hits@10} & \textbf{Mean Rank} \\
\midrule
TransE & 0.244 & 0.110 & 0.301 & 0.524 & 32.00 \\
ComplEx & 0.027 & 0.004 & 0.008 & 0.041 & 100.81 \\
\bottomrule
\end{tabular}
\caption{Validation link prediction metrics for TransE and ComplEx}
\end{table}

\subsection{Final Test Evaluation (TransE)}
The best performing model on the validation set, TransE, was evaluated on the test set:
\begin{itemize}
    \item \textbf{Filtered MRR:} 0.240
    \item \textbf{Hits@1:} 0.101
    \item \textbf{Hits@3:} 0.286
    \item \textbf{Hits@10:} 0.532
    \item \textbf{Mean Rank:} 30.74
\end{itemize}

A stability check across three random seeds (7, 42, 1337) yielded test MRR values of 0.256, 0.240, and 0.296, indicating stable learning with a mean MRR of approximately 0.264 (SD = 0.029).

\section{Interpretation}
TransE had a considerably higher filtered MRR (0.244) than ComplEx (0.027) on this dataset and training budget. 
The best model (TransE) is moderately effective at placing held-out true links near the top of the candidate list, as shown by its Hits@10 of 0.524, meaning that over half of the time the true missing link was in the top 10 candidates.
The global filtered MRR measures the model's ability to rank withheld true triples among candidate entity replacements across the whole graph. It is not alert precision, attack recall or a calibrated likelihood of compromise.

\section{Addressing Learning Outcomes (LO1--LO12)}

\textbf{LO1 (Knowledge Graph Embeddings):} TransE represents relations as translations in embedding space; ComplEx uses complex-valued embeddings and can represent more varied relational patterns. The experiment compares them through the same held-out link prediction task. While ComplEx has the capacity for more complex patterns, TransE was more easily optimized within the limited training budget of 80 epochs on this relatively simple graph.

\textbf{LO2 (Logical knowledge in KGs):} The rule mapping \texttt{has\_family} to \texttt{suggests\_technique} is explicit, inspectable and separate from training. This shows why a symbolic layer can produce provenance-aware facts. The code does not implement full recursive or existential Datalog reasoning, but it demonstrates the integration of simple deductive logic into a KG pipeline.

\textbf{LO3 (Graph Neural Networks):} While not explicitly implemented in this baseline feasibility study, GNNs could serve as a powerful alternative to translational KGEs. By message-passing over the observed \texttt{runs} or \texttt{connects\_to} subgraph structure, a GNN might directly capture structural neighborhoods before predicting \texttt{suggests\_technique} edges.

\textbf{LO4 (KG Data Models):} The generator naturally resembles a property graph because log events are typed entities with labelled edges. PyKEEN consumes three-column labelled triples, which resembles an RDF-style statement view. This is a modelling choice, not proof that either model is universally better.

\textbf{LO5 (KG Architectures):} The architecture is modular: ingestion/generation $\rightarrow$ rule engine $\rightarrow$ triple store $\rightarrow$ KGE training/evaluation $\rightarrow$ ranked analyst output. In a real system, raw telemetry would live outside the embedding model and only selected, versioned KG facts would be embedded.

\textbf{LO6 (Scalable Reasoning):} The KGE serves as approximate KG completion by ranking missing links. The rule engine serves as deterministic reasoning. These are different forms of inference and should not be conflated.

\textbf{LO7 (KG Creation):} The knowledge graph was created via an ingestion pipeline that generated a schema of typed entities from pseudonymous logs, mimicking record linkage and schema mapping by standardizing edge types like \texttt{executed\_by}.

\textbf{LO8 (KG Evolution):} The KGE pipeline represents Knowledge Graph completion (adding to it via link prediction). Predicting a missing \texttt{suggests\_technique} link evolves the graph with new probable hypotheses. 

\textbf{LO9 (Real-world Applications):} A SOC analyst could review a ranked list of candidate technique links alongside the rule provenance and raw evidence. Human review is necessary because KGE scores are not calibrated security risk values.

\textbf{LO10 (Financial KG Applications):} While this dataset simulates IT security events (which are highly relevant to banks protecting infrastructure from hostile attacks), in financial institutions, knowledge graphs also map company ownership and regulatory constraints. Similar link prediction could help uncover hidden relations in hostile takeover scenarios.

\textbf{LO11 (KG Services):} The ranked technique-tail metrics simulate a query interface for analysts: "Which ATT\&CK techniques are most likely given this process?" By retrieving these lists, the KG provides an analytical service powered by deep KGE-based knowledge.

\textbf{LO12 (Connections between KGs, ML and AI):} The project demonstrates a hybrid AI design: deductive logic creates transparent facts, while ML generalises patterns in the graph. Neither component alone makes the system a complete AI-based security product, but their combination forms a "melting pot" characteristic of modern KG-AI architectures.

\section{Limitations}
\begin{enumerate}
    \item The data are synthetic and intentionally regular, so results cannot be generalised to real APT detection.
    \item The random triple split is appropriate for a teaching KGE benchmark but can leak temporal information in a log setting. A real extension needs a chronological split.
    \item Ranking metrics test the ability to recover withheld triples, not alert precision/recall on independently labelled incidents.
    \item The rules are intentionally simple and do not cover MITRE ATT\&CK in full.
    \item The target technique links are derived from a deterministic family-to-technique rule. Consequently, the experiment evaluates whether a KGE can recover a structured, rule-generated pattern from the rest of the graph; it does not show that the model detects independently labelled APT activity.
    \item Hyperparameters are small and illustrative; no comprehensive HPO was run.
\end{enumerate}

\section{Safe Next Extension}
Replace \texttt{generate\_base\_triples()} with a converter for a small, documented window of one data source. Keep stable pseudonymous identifiers, add timestamp handling separately, create a chronological split, and retain the exact rule set and configuration used for each run.

\end{document}
